%% Dans ce fichier, les �ventuels accents sont cod�s
%% en ISO-8859-15 (aka latin9)
%% La ligne suivante donne trois exemples 
%% e accent aigu : �, e accent grave : �, e-dans-l'o : �
%%%%%ENTREES%%%%
\begin{tabarg}
\multicolumn{3}{l}{\titretab {Entr\'ees}}\\
\hline
$X$ & Vecteur, \texttt{float} & Param�tre X \\
$Y$ & Vecteur, \texttt{float} & Param�tre Y \\
$f_e$ & Coeff., \texttt{float} & Fr�quence d'�chantillonnage X et Y (Hz)\\
$T_1$ & Coeff., \texttt{int} & Dur�e de la p�riode d'int�gration partielle  (s)\\
$T_2$ & Coeff., \texttt{int} & Dur�e de la p�riode d'int�gration totale  (s)\\
$R_{manq}$ & Coeff., \texttt{int} & Fraction minimale de valeurs non manquantes tol�r�es (compris entre 0 et 1) \\
\end{tabarg}

\medskip
%%%%%SORTIES%%%%
\begin{tabarg}
\multicolumn{3}{l}{\titretab {Sorties}}\\
\hline
$I_{stat}$ & Vecteur, \texttt{float} & Indice de stationnarit� (sans unit�)\\
\end{tabarg}
